\input{../tex-inputs/latex-leseansicht-vorspann}

\section[Gustav Schwarzkopf an Arthur Schnitzler, 4. 7. 1913]{L04259 Gustav Schwarzkopf an Arthur Schnitzler, 4. 7. 1913}
\nopagebreak\mylabel{L04259v}
\rehead{ }\normalsize\beginnumbering\briefempfaengerindex{Schnitzler, Arthur@\textsc{Schnitzler, Arthur}!zzzSchwarzkopf, Gustav@\emph{von Gustav Schwarzkopf}!1913-07-041@{4. 7. 1913}|(be}
\toendnotes[C]{\smallbreak\pagebreak[2]}
\correspDesc{Versand  durch Gustav Schwarzkopf am 4. 7. 1913 in Opatija
\newline{}Erhalt  durch Arthur Schnitzler im Zeitraum [8. 4. 1913 – 12. 4. 1913?] in Wien}\toendnotes[C]{\smallbreak}
\Standort{CUL, Schnitzler, B 96a.}
\physDesc{Postkarte, 638 Zeichen
\newline{}Handschrift: schwarze Tinte, deutsche Kurrent
\newline{}Versand: Stempel: »\nobreak{}\oindex{Opatija@\textbf{Opatija}|pwk}Abbazia, 4 VII 13, \textcolor{gray}{7}\nobreak{}«.  
\newline{}Schnitzler: mit Bleistift Vermerk: »\textsc{Schwarzkopf}« }\toendnotes[C]{\smallbreak}\pstart{}{\pb}\textsc{Abbazia. Wiener Heim\oindex{Pension Wiener Heim@\textbf{Pension Wiener Heim}|pw}}\pend{}{\bigskip}\pstart{}Herrn\pend{}\pstart{}\textsc{D\textsuperscript{r} Arthur Schnitzler}\pend{}\pstart{}\textsc{Wien XVIII\oindex{XVIII., Währing@\textbf{XVIII., Währing}|pw}}\pend{}\pstart{}Sternwarteſtraße 71\oindex{Wien@\textbf{Wien}!XVIII., Währing@\textbf{XVIII., Währing}!Sternwartestraße 71@\textbf{Sternwartestraße 71}|pw}\pend{}{\bigskip}\vspace{1em}
\pstart
           \raggedleft{}{\pb}Abbazia\oindex{Opatija@\textbf{Opatija}|pw}{ }4. VII 13.\pend
           
\pstart
           \raggedleft{}Wiener Heim\oindex{Pension Wiener Heim@\textbf{Pension Wiener Heim}|pw}, Zim 19.\pend
           \vspace{0.5em}
\pstart
           Lieber Arthur, ich bin heute Vormittag hier angeko{\geminationm}en, unter Blitz und Do{\geminationn}er,
               im{ }ſtrömenden Regen, und jetzt – 3 Uhr —{ }ſchüttet es noch i{\geminationm}er. Von den Reizen der Natur habe ich also noch nichts
               bemerken kö{\geminationn}en und der Tag unterſcheidet{ }ſich bis jetzt
               nicht viel von meinem letzten Tag in Wien\oindex{Wien@\textbf{Wien}|pw}, wo es
               auch unaufhörlich geregnet hat. Übrigens{ }ſcheine ich ganz gut untergebracht zu{ }ſein –
               nach dem erſten Mittageſſen zu{ }ſchließen. Auch \textsc{\label{K_L04259-1v}\edtext{Laſsnitz\oindex{Laßnitzhöhe@\textbf{Laßnitzhöhe}|pw}}{\lemma{\textnormal{\emph{Lassnitz}}}\Cendnote{\textnormal{Obwohl es mehrere
                     Orte mit diesem Namen gibt, dürfte der zu dieser Zeit als Kurort im Aufschwung befindliche Ort Lassnitzhöhe\oindex{Laßnitzhöhe@\textbf{Laßnitzhöhe}|pwk} 15km östlich von Graz\oindex{Graz@\textbf{Graz}|pwk} gemeint sein.
                     Die dazugehörige Bahnstation  Laßnitzthal\oindex{Bahnhof Laßnitzthal@\textbf{Bahnhof Laßnitzthal}|pwk} hieß damals nur Laßnitz\oindex{Bahnhof Laßnitzthal@\textbf{Bahnhof Laßnitzthal}|pwk}.
                  }}}\label{K_L04259-1}} war ganz nett, da war das Wetter noch gnädiger.\pend
           
\pstart
           Ihre Frau\pwindex{Schnitzler, Olga 17.\,1.\,1882 Wien – 13.\,1.\,1970 Lugano@\textsc{Schnitzler, Olga} (17.\,1.\,1882 Wien – 13.\,1.\,1970 Lugano)|pwv} und Sie herzlichſt grüßend{\\[\baselineskip]}Ihr{\\[\baselineskip]}\spacefill\mbox{Gustav.}\pend
           \leftskip=0em{}\selectlanguage{ngerman}\endnumbering\briefempfaengerindex{Schnitzler, Arthur@\textsc{Schnitzler, Arthur}!zzzSchwarzkopf, Gustav@\emph{von Gustav Schwarzkopf}!1913-07-041@{4. 7. 1913}|)be}\mylabel{L04259h}
\begin{anhang}
\end{anhang}\newcommand{\dateiname}{L04259}\newcommand{\titel}{Gustav Schwarzkopf an Arthur Schnitzler, 4. 7. 1913}\newcommand{\editorInnen}{Herausgegeben von Jahnke, SelmaMüller, Martin Anton}\input{../tex-inputs/latex-leseansicht-abspann}
